% !TEX root = ../main.tex
\begin{ouverture}[trigo]{6}{Trigonométrie}{Partie 2 · Analyse}
{Radian\enspace·\enspace Cercle unité\enspace·\enspace Sinus et cosinus\enspace·\enspace Angles associés}
\objectif{Passer des degrés aux radians et calculer une longueur d’arc.}
\objectif{Placer un point sur le cercle trigonométrique.}
\objectif{Déterminer des valeurs exactes par symétrie.}
\objectif{Utiliser les coordonnées du cercle dans une situation géométrique.}
\prerequis{Triangles rectangles et théorème de Pythagore.}
\prerequis{Coordonnées et symétries dans le plan.}
\end{ouverture}
\essentiel
\essSection{1}{Enrouler une droite sur le cercle}
\begin{definition}
Le cercle trigonométrique est le cercle de centre $O$ et de rayon 1, orienté
dans le sens inverse des aiguilles d'une montre. Un réel $t$ correspond à un
déplacement de longueur $|t|$ sur ce cercle : sens positif si $t>0$, sens
négatif si $t<0$. Un tour correspond à $2\pi$ radians, soit $360^\circ$.
Ainsi $180^\circ=\pi$ rad et $90^\circ=\pi/2$ rad.
\end{definition}
\begin{propriete}
Sur un cercle de rayon $R$, un angle de mesure $\theta\geq0$ radians intercepte
un arc de longueur $R\theta$. Deux réels ont le même point-image exactement
lorsque leur différence est un multiple entier de $2\pi$.
\end{propriete}
\essSection{2}{Les coordonnées du point-image}
\begin{definition}
Le point-image de $t$ a pour coordonnées $(\cos t;\sin t)$ dans le repère
orthonormé $(O;I;J)$. On a $-1\leq\cos t\leq1$, $-1\leq\sin t\leq1$ et
$\cos^2t+\sin^2t=1$.
\end{definition}
\begin{center}
\begin{tikzpicture}[scale=1.45]
\draw[otMuted,->](-1.45,0)--(1.6,0)node[right]{$\cos t$};
\draw[otMuted,->](0,-1.25)--(0,1.45)node[above]{$\sin t$};
\draw[otNavy,thick](0,0)circle(1);
\draw[otGreen,thick](0,0)--(.5,.8660254);
\draw[otGreenDark,->](.3,0)arc(0:60:.3);
\draw[otMuted,dashed](.5,0)--(.5,.8660254)--(0,.8660254);
\fill[otGreenDark](.5,.8660254)circle(1.4pt)node[above right]{$M$};
\node[below left]at(0,0){$O$};\node[below right]at(1,0){$I$};
\node[font=\small]at(.45,.22){$t$};
\end{tikzpicture}
\par\legende{Pour $t=\pi/3$, le point $M$ a pour coordonnées $(1/2;\sqrt3/2)$.}
\end{center}

\clearpage
\essSection{3}{Les valeurs remarquables}
{\renewcommand{\arraystretch}{1.9}\begin{tabularx}{\linewidth}{@{}l*{5}{>{\centering\arraybackslash}X}@{}}
\toprule $t$ & $0$ & $\pi/6$ & $\pi/4$ & $\pi/3$ & $\pi/2$\\\midrule
$\cos t$ & $1$ & $\sqrt3/2$ & $\sqrt2/2$ & $1/2$ & $0$\\
$\sin t$ & $0$ & $1/2$ & $\sqrt2/2$ & $\sqrt3/2$ & $1$\\\bottomrule
\end{tabularx}}
\begin{propriete}[Les symétries du cercle]
Pour tout réel $t$,
\[\cos(-t)=\cos t,\quad\sin(-t)=-\sin t,\]
\[\cos(\pi-t)=-\cos t,\quad\sin(\pi-t)=\sin t,\]
\[\cos(t+\pi)=-\cos t,\quad\sin(t+\pi)=-\sin t.\]
Ajouter $2k\pi$, avec $k$ entier relatif, ne change ni sinus ni cosinus.
\end{propriete}
\begin{savoirfaire}[Déterminer une valeur exacte]
Réduire l'angle d'un nombre entier de tours. Repérer le quadrant. Utiliser
une symétrie pour revenir à une valeur remarquable, puis vérifier le signe.
Par exemple, $\cos(5\pi/6)=-\cos(\pi/6)=-\sqrt3/2$.
\end{savoirfaire}
\begin{attention}
Les mesures utilisées dans les calculs sont en radians. Un angle peut avoir
plusieurs mesures, mais un point-image est unique. Une valeur décimale à la
calculatrice ne remplace pas une valeur exacte demandée.
\end{attention}
\begin{exemple}[Un argument géométrique]
Dans un triangle rectangle isocèle de côtés de l'angle droit égaux à 1,
l'hypoténuse vaut $\sqrt2$ ; donc $\cos(\pi/4)=\sin(\pi/4)=1/\sqrt2$.
La démonstration des autres valeurs est détaillée dans les démonstrations de référence.
\end{exemple}

\exercices

\rubrique{Découverte}
\begin{blocLivre}{ex-6-1}
\begin{exercice}[1]{Convertir des angles}
\begin{questions}
\item Convertir $30^\circ$, $45^\circ$, $120^\circ$ et $270^\circ$ en radians.
\item Convertir $\pi/3$, $5\pi/6$ et $7\pi/4$ en degrés.
\end{questions}
\end{exercice}
\end{blocLivre}
\begin{blocLivre}{ex-6-2}
\begin{exercice}[1]{Même point-image}
\begin{questions}
\item Les réels $\pi/3$ et $13\pi/3$ ont-ils le même point-image ?
\item Même question pour $-\pi/2$ et $3\pi/2$, puis pour $\pi/4$ et $5\pi/4$.
\end{questions}
\label{figlivre-r06-02-images-ex}
\figureLivre{r06-02-images}

\end{exercice}
\end{blocLivre}
\begin{blocLivre}{ex-6-3}
\begin{exercice}[1]{Lire les valeurs remarquables}
\begin{questions}
\item Donner les coordonnées des points-images de $0$, $\pi/6$, $\pi/4$ et $\pi/2$.
\item Donner $\cos\pi$, $\sin\pi$, $\cos(3\pi/2)$ et $\sin(3\pi/2)$.
\end{questions}
\label{figlivre-r06-03-valeurs-ex}
\figureLivre{r06-03-valeurs}

\end{exercice}
\end{blocLivre}
\begin{blocLivre}{ex-6-4}
\begin{exercice}[1]{Utiliser une symétrie}
\begin{questions}
\item Calculer exactement $\cos(-\pi/3)$ et $\sin(-\pi/3)$.
\item Calculer $\cos(2\pi/3)$ et $\sin(2\pi/3)$.
\end{questions}
\label{figlivre-r06-04-symetries-ex}
\figureLivre{r06-04-symetries}

\end{exercice}
\end{blocLivre}
\begin{blocLivre}{ex-6-5}
\begin{exercice}[1]{Une longueur d'arc}
Un cercle a pour rayon 6 cm. Calculer la longueur des arcs associés aux angles $\pi/3$ et $3\pi/4$ radians, puis la longueur du cercle entier.
\label{figlivre-r06-05-arcs-ex}
\figureLivre{r06-05-arcs}

\end{exercice}
\end{blocLivre}
\clearpage
\rubrique{Entraînement}
\begin{blocLivre}{ex-6-6}
\begin{exercice}[2]{Réduire des grandes mesures}
\begin{questions}
\item Ramener $17\pi/6$, $-11\pi/4$ et $29\pi/3$ dans $\mathopen{]}-\pi;\pi]$.
\item En déduire le sinus et le cosinus de chacune.
\end{questions}
\end{exercice}
\end{blocLivre}
\begin{blocLivre}{ex-6-7}
\begin{exercice}[2]{Retrouver une coordonnée}
On sait que $\cos t=-3/5$ et que $t\in[\pi/2;\pi]$.\begin{questions}
\item Calculer $\sin^2t$ puis $\sin t$.
\item Déterminer les coordonnées des points-images de $-t$ et $t+\pi$.
\end{questions}
\label{figlivre-r06-07-coordonnee-ex}
\figureLivre{r06-07-coordonnee}

\end{exercice}
\end{blocLivre}
\begin{blocLivre}{ex-6-8}
\begin{exercice}[2]{Deux angles associés}
\begin{questions}
\item Calculer $\cos(5\pi/4)$, $\sin(5\pi/4)$, $\cos(7\pi/6)$ et $\sin(7\pi/6)$.
\item Calculer $A=2\cos(5\pi/4)+\sqrt2\sin(7\pi/6)$.
\end{questions}
\end{exercice}
\end{blocLivre}
\begin{blocLivre}{ex-6-9}
\begin{exercice}[2]{Retrouver des angles sur le cercle}
À l’aide du cercle, déterminer tous les réels $t\in[0;2\pi[$ tels que :\begin{questions}
\item $\cos t=1/2$.
\item $\sin t=-\sqrt2/2$.
\item $\cos t=0$.
\end{questions}

\label{figgeo-06-09-ex}
\figureGeo{06-09}
\end{exercice}
\end{blocLivre}
\begin{blocLivre}{ex-6-10}
\begin{exercice}[2]{Une portion de cercle}
Un disque de rayon 10 cm est partagé en un secteur d’angle $2\pi/5$ et son complément.\begin{questions}
\item Calculer la fraction du disque occupée par le premier secteur.
\item Calculer son aire et son périmètre, en incluant les deux rayons.
\end{questions}

\label{figgeo-06-10-ex}
\figureGeo{06-10}
\end{exercice}
\end{blocLivre}
\begin{blocLivre}{ex-6-11}
\begin{exercice}[2]{Une identité à démontrer}
\begin{questions}
\item Développer $(\cos t+\sin t)^2+(\cos t-\sin t)^2$.
\item En déduire sa valeur pour tout réel $t$.
\item Vérifier le résultat pour $t=\pi/4$ sans calculatrice.
\end{questions}
\end{exercice}
\end{blocLivre}
\clearpage
\rubrique{Problèmes}
\begin{blocLivre}{ex-6-12}
\begin{exercice}[3]{Une grande roue}
Une roue tourne à vitesse de rotation uniforme. Son rayon est de 12 m et son centre est à 14 m du sol. Le point $M$ a pour angle $t$ à partir de l’horizontale positive. Ses coordonnées dans un repère au sol sont $M(12\cos t;14+12\sin t)$.\begin{questions}
\item Donner les positions de $M$ pour $t=0$, $\pi/2$, $\pi$ et $3\pi/2$.
\item Calculer la hauteur pour $t=\pi/6$ et $7\pi/6$.
\item Si un tour dure 48 s, quelle durée correspond à un angle parcouru de $5\pi/6$ ?
\end{questions}

\label{figgeo-06-12-ex}
\figureGeo{06-12}
\end{exercice}
\end{blocLivre}
\begin{blocLivre}{ex-6-13}
\begin{exercice}[3]{Un hexagone régulier}
On inscrit un hexagone régulier dans le cercle unité, un sommet étant $A(1;0)$.\begin{questions}
\item Donner les angles puis les coordonnées de ses six sommets.
\item Calculer la longueur d’un côté avec la distance entre les deux premiers sommets.
\item En déduire le périmètre et l’aire exacte de l’hexagone.
\end{questions}

\label{figgeo-06-13-ex}
\figureGeo{06-13}
\end{exercice}
\end{blocLivre}
\begin{blocLivre}{ex-6-14}
\begin{exercice}[3]{Une zone autorisée}
Un point $M(\cos t;\sin t)$ appartient au cercle unité, avec $t\in[0;2\pi[$.\begin{questions}
\item Représenter la zone où $\sin t\geq0$.
\item Déterminer les angles où $\cos t\leq-1/2$.
\item Déterminer les angles qui satisfont simultanément les deux conditions et calculer la longueur de l’arc correspondant.
\end{questions}

\label{figgeo-06-14-ex}
\figureGeo{06-14}
\end{exercice}
\end{blocLivre}
\begin{blocLivre}{ex-6-15}
\begin{exercice}[3]{Approcher un cercle par des polygones}
Un polygone régulier à $n\geq3$ côtés est inscrit dans le cercle unité. En partageant un triangle central en deux, on obtient un demi-côté de longueur $\sin(\pi/n)$.\begin{questions}
\item Exprimer le périmètre $P_n$ de ce polygone.
\item Calculer exactement $P_6$ et $P_{12}$ en utilisant $\sin(\pi/12)=(\sqrt6-\sqrt2)/4$, fourni ici.
\item Comparer numériquement ces périmètres à $2\pi$.
\item Méthode d’Archimède. On admet que le côté $c_{2n}$ du polygone à $2n$ côtés vérifie $c_{2n}=\sqrt{2-\sqrt{4-c_n^2}}$. En partant de l’hexagone, de côté $c_6=1$, écrire une fonction Python, sans utiliser $\pi$, qui renvoie le demi-périmètre du polygone à $6\times2^k$ côtés. Donner une valeur approchée pour 96 côtés.
\end{questions}

\label{figgeo-06-15-ex}
\figureGeo{06-15}
\end{exercice}
\end{blocLivre}
